\section{Pre-v3 Generalization Studies}
Before the modern freeze, the controller family was examined across datasets
and association backends. These studies are kept separate from v3 because
their detectors, evaluators, development splits, and action definitions are
not identical.

\subsection{Historical UAVDT transfer}
The frozen historical profiles were transferred to UAVDT without retuning. The
baseline obtained HOTA/MOTA/IDF1/IDS $24.085/13.841/27.887/558$; V1 obtained
$28.390/17.399/34.453/321$; and V2 obtained $26.930/16.118/32.014/308$.
This supports portability of frozen operating profiles, not a universal
quality ordering and not evidence that the modern v3 policy was trained on
UAVDT.

\subsection{U2MOT host on VisDrone}
The cross-pipeline study used the published U2MOT host with YOLOX-X
\cite{liu2023u2mot,ge2021yolox}. Development used seven VisDrone validation
sequences and held-out evaluation used 17 VisDrone test-dev sequences. The
V1 SCI weights and transforms were retained, the tracker was not modified,
and the action boundaries were recalibrated from the U2MOT validation
distribution using unsupervised 33rd/67th percentiles because the V1
thresholds placed all frames in the hard region. No ground-truth metric was
used for that recalibration, but the U2MOT action tuples were hand-set and do
not have a selection record.

On test-dev, the custom evaluator reported baseline MOTA 53.9, IDF1 69.8,
IDS 1239, FP 41385, and FN 63241; the transferred controller reported MOTA
53.7668, IDF1 69.8494, IDS 1152, FP 40155, and FN 64801. The defensible
summary is approximately unchanged accuracy, 87 fewer IDS, 1230 fewer FP, and
1560 more FN. No HOTA, confidence interval, or matched-baseline speed claim
is attached. This is a VisDrone cross-pipeline study, not a UAVDT result or a
universal quality win.

\subsection{SparseTrack host on MOT17}
The separate SparseTrack study used the published pseudo-depth association
host \cite{liu2025sparsetrack} on MOT17 \texttt{val\_half}, seven FRCNN
sequences, and 2652 frames. Its minimal Adaptive Edge V1 transfer changed only
NMS from 0.70 to 0.80; confidence, resolution, detector, and tracker settings
were unchanged. The threshold came from a leave-one-sequence-out stump study
on those development sequences. Static versus adaptive evaluation gave MOTA
76.8881 versus 76.9252, IDF1 81.4943 versus 81.5299, and IDS 136 versus 123.
The MOTA change was only $+0.0371$ percentage points with bootstrap CI
$[-0.041,+0.182]$, so the manifest does not establish superiority. This
near-neutral result is retained because it limits claims of broad
cross-dataset generalization.
